\documentclass{amsart}
\usepackage{amsmath, amssymb}
\usepackage{physics}

% For \includegraphics 
\usepackage{graphicx}

\title{Week 002}
\author{Kazuya Haine}
\date{}

\newcommand{\Z}{\mathbb{Z}}
\newcommand{\C}{\mathbb{C}}
\newcommand{\calS}{\mathcal{S}}

\newtheorem*{definition}{Definition}
\newtheorem*{remark}{Remark}

\newtheorem{theorem}{Theorem}

\newcounter{qpart}
\newcommand{\qpart}[1]{%
  \stepcounter{qpart}%
  \bigskip\noindent\textbf{(\alph{qpart})}\quad\textit{#1}\par\bigskip%
}

\begin{document}
\maketitle

\begin{definition}[Additive uniqueness]
    Let $\calS$ be a subset of the integers, $\Z$. $\calS$ is said to satisfy \underline{additive uniqueness} if, whenever $n_k \in \calS$ for $k = 1,2,3,4$ and $n_1 + n_2 = n_3 + n_4$, then either $n_1 = n_3$ or $n_1 = n_4$.
\end{definition}
\begin{remark}
    This can be interpreted as saying that addition is injective on $\calS$ when seen as an operation on unordered pairs.
\end{remark}
\begin{center}
\rule{0.9\textwidth}{0.5pt}
\end{center}

\bigskip
\noindent
Let $\calS$ be a subset of $\Z$ with $N \ge 1$ elements satisfying additive uniqueness. Define $f \colon \C \to \C$ by
\begin{equation*}
    f(x) := \sum_{n\in\calS} e^{2\pi inx}.
\end{equation*}

\qpart{Evaluate the following: \begin{equation*} (1) \quad \int_0^1 |f(x)|^2 \dd{x} \qquad (2) \quad \int_0^1 |f(x)|^4 \dd{x}\end{equation*}}

\noindent
$|f(x)|^2$ can be expressed as the product of $f(x)$ with it's complex conjugate:
\begin{align*}
    |f(x)|^2 &= f(x) \overline{f(x)} \\
             &= \left(\sum_{m \in \calS} e^{2\pi imx}\right)\left(\sum_{n \in \calS} e^{-2\pi inx}\right) \\
             &= \sum_{m \in \calS} \sum_{n \in \calS} e^{2\pi i(m-n)x} \\
             &= \left(\sum_{n \in \calS} 1\right) + \left(\sum_{\substack{n,m \in \calS, \\ n < m}} 4\times\frac{e^{i(2\pi(m-n)x)} + e^{-i(2\pi(m-n)x)}}{2}\right) \\
             &= N + \sum_{\substack{n,m \in \calS, \\ n < m}} 4\cos(2\pi(m-n)x)
\end{align*}
Thus, 
\begin{align*}
    \int_0^1 |f(x)|^2 \dd{x} &= \int_0^1 \left(N + \sum_{\substack{n,m \in \calS, \\ n < m}} 4\cos(2\pi(m-n)x) \right)\dd{x} \\
                             &= N,
\end{align*}
since every cosine term contributes 0 to the integral.

\noindent
We can use a similar trick to evaluate the second integral by considering $|f(x)|^4$ as $f(x)^2 \overline{f(x)^2}$. We first expand $f(x)^2$:
\begin{align*}
    f(x)^2 &= \left(\sum_{m \in \calS} e^{2\pi imx}\right)\left(\sum_{m \in \calS} e^{2\pi imx}\right) \\
           &= \sum_{m \in \calS} \sum_{n \in \calS} e^{2\pi i(m+n)x} \\
           &= \sum_{n \in \calS} e^{2\pi i(2n)x} + 2\sum_{\substack{n,m \in \calS \\ n < m}} e^{2\pi i(n + m)x}
\end{align*}
Here, additive uniqueness tells us that each of the terms $2n$, and $m + n$ indexed in these ways are distinct. thus, $f(x)^2$ essentially takes the form of $f(x) + 2g(x)$ where $g$ is the corresponding function to $f$ when replacing $\calS$ with $\{ n + m \mid n,m \in \calS \}$.

\bigskip\noindent
Repeating the expansion of the absolute square as the product with the complex conjugate:
\begin{align*}
    |f(x)|^4 &= f(x)^2 \overline{f(x)^2} \\
             &= (f(x) + 2g(x))\overline{(f(x) + 2g(x))} \\
             &= f(x)\overline{f(x)} + 4g(x)\overline{g(x)} + 2\bigg(f(x)\overline{g(x)} + g(x)\overline{f(x)}\bigg)
\end{align*}
We know how $f(x)\overline{f(x)}$ expands and what it contributes to the integral: $N$, and since $g(x)\overline{g(x)}$ has the same structure, it contributes $\# \{n + m \mid n, m \in \calS \} = \frac{N(N-1)}{2}$.
Finally, the $f(x)\overline{g(x)} + g(x)\overline{f(x)}$ corresponds to more cosine terms with non-zero integer multiples of $2\pi$ periods in the expansion, so it contributes 0.
So, together,
\begin{equation*}
    \int_0^1 |f(x)|^4 \dd{x} = N + 2N(N-1) = 2N^2 - N.
\end{equation*}

\qpart{Prove that \begin{equation*} \int_0^1 |f(x)| \dd{x} \ge \frac{1}{2\sqrt{N}}.\end{equation*} When does the equality hold?}

\noindent
By Cauchy-Schwarz, we have,
\begin{align*}
    \left(\int_0^1 |f(x)|^2 \dd{x}\right)^2 &\le \left(\int_0^1 |f(x)| \dd{x}\right) \left(\int_0^1 |f(x)|^3 \dd{x}\right) \\
    N^2                                     &\le \left(\int_0^1 |f(x)| \dd{x}\right) \left(\int_0^1 |f(x)|^3 \dd{x}\right).
\end{align*}
A second application of Cauchy-Schwarz, gives:
\begin{align*}
    \left(\int_0^1 |f(x)|^3 \dd{x}\right)^2 &\le \left(\int_0^1 |f(x)|^2 \dd{x}\right) \left(\int_0^1 |f(x)|^4 \dd{x}\right) \\
    \left(\int_0^1 |f(x)|^3 \dd{x}\right)^2 &\le N^2(2N-1) \\
    \int_0^1 |f(x)|^3 \dd{x} &\le N\sqrt{2N - 1}.
\end{align*}
Putting them together, we have:
\begin{equation*}
    \int_0^1 |f(x)| \dd{x} \ge \frac{N^2}{N\sqrt{2N-1}} = \frac{N}{\sqrt{2N-1}}
\end{equation*}
For all $N \ge 1$, $\frac{N}{\sqrt{2N-1}} > \frac{1}{2\sqrt{N}}$ since rearranging this gives the equivalent inequality $4N^3 > 2N - 1$ which is clearly true for all such $N$, so
\begin{equation*}
    \int_0^1 |f(x)| \dd{x} \ge \frac{N}{\sqrt{2N-1}} > \frac{1}{2\sqrt{N}}
\end{equation*}
\bigskip\noindent
This tells us that the equality never holds.
\end{document}
